\section{Restaurant Customers}
\label{sec:g-restaurant}
\source{CSES 1619 (Notion: Greedy)}{Easy}

\problemstatement{Customer $i$ is in the restaurant from time $a_i$ to $b_i$
($1\le a_i<b_i\le10^9$, all $2n$ times distinct, $n\le2\cdot10^5$). Find the
maximum number of customers present at the same time.}

\begin{patternbox}
``Maximum overlap of intervals'' $\Rightarrow$ \textbf{sweep line}: an
arrival is a $+1$ event, a departure a $-1$ event; the answer is the maximum
prefix sum in time order.
\end{patternbox}

\subsection*{Approach and intuition}

The number of people present only changes at event times. Sorting the $2n$
events and accumulating $\pm1$ reproduces the head count at every moment; its
maximum is the answer. Because times are distinct we never have to decide
whether a departure and an arrival at the same instant overlap.

\subsection*{Algorithm}
\begin{algorithmic}[1]
\State events $\gets \{(a_i,+1),(b_i,-1)\}$, sort by time
\State $cur\gets0$, $best\gets0$
\For{each event $(t,d)$} \State $cur\gets cur+d$, $best\gets\max(best,cur)$ \EndFor
\State \Return $best$
\end{algorithmic}

\dryrun{Intervals $(5,8),(2,4),(3,9)$: events $2^{+},3^{+},4^{-},5^{+},8^{-},9^{-}$,
counts $1,2,1,2,1,0$. Maximum $2$. \checkmark}

\subsection*{C++17 implementation}
\cppfile{greedy/08_restaurant_customers.cpp}

\begin{complexitybox}
\textbf{Time:} $O(n\log n)$. \textbf{Space:} $O(n)$.
\end{complexitybox}

\begin{pitfallbox}
\begin{itemize}
  \item If times could coincide you must decide the tie order (process
        $-1$ before $+1$ for half-open intervals). Here they are distinct.
  \item Coordinates up to $10^9$: never allocate an array indexed by time.
\end{itemize}
\end{pitfallbox}
